\chapter{Összefoglalás}

\section{Feladat}
\paragraph{}A dolgozat célja az volt, hogy egy cukorbetegség kezelését támogató alkalmazás 
étkezésnaplózó funkciójához kereskedelmi API-n elérhető multimodális nagy nyelvi modellek 
pontosságát és használhatóságát értékelje, és megvizsgálja, hogy a becslés javítható-e. Ehhez öt 
képes és egy szöveges adathalmazon egységes mérési keretrendszert készítettem, amely a modellek 
válaszát JSON-séma szerint kéri, a becsléseket a referenciaértékekkel veti össze, és minden futás 
beállításait rögzíti. Az 50 ételes előmérések eredményei alapján előre rögzített döntési mérést 
terveztem, amelynek hipotéziseit az előmérésben fel nem használt ételeken, Holm-korrekcióval 
vizsgáltam. Az ételek és összetevők megnevezését leíró jelleggel értékeltem. A dolgozat a modellek és 
a feldolgozási stratégiák értékelésére összpontosít, ezért a témabejelentőben szereplő integrált 
implementáció helyett a mérések alapján javasolt becslési folyamatot kitekintésként \az{\ref{sec:javasolt_folyamat}}.~alfejezet írja le.

\section{Eredmények}
\paragraph{}A modellek egymáshoz viszonyított pontossága adathalmazonként eltért. A döntési mérésben a Flash-Lite szénhidrát-pMAE-je 23,2 és 53,6\% között volt. Az 
adathalmazok konyhája, fényképezési körülményei és referenciaértékeinek forrása is eltér, 
és a referenciaértékek szóródása is különböző, ezért a pMAE az adathalmazok között nem hasonlítható 
össze. A DFood-Tj-n például a referenciaértékek kis szóródása miatt már a mediánbecslés is kis pMAE-t ad. 
A Flash-Lite a Nutrition5k-n, az OmniFood8k-n és az ACETADA-n szénhidrátban szignifikánsan pontosabb volt 
a nála nagyjából öt-hatszor drágább Sonnetnél és Solnál, a DFood-Tj-n viszont mindkét drágább modell volt 
a pontosabb. Ez utóbbi eredmény nem választható el attól, hogy a DFood-Tj képeinek egy részén a 
mérleg kijelzője is leolvasható. Más célmennyiségekben a sorrend eltért, így egyik modell sem volt 
minden adathalmazon és minden célmennyiségben a legpontosabb.

\paragraph{}Az előmérés hibaforrás-bontása szerint az OmniFood8k-n az energiahiba szinte teljesen az 
energiasűrűség becsléséből adódott, a SNAPMe-n inkább a tömegéből, a többi adathalmazon a két rész 
hasonló nagyságú volt, és a Flash-Lite a nagy adagokat a kicsikhez képest mindenhol alulbecsülte. 
Több kép megadása az előmérésben nem javított egyértelműen, és a promptváltozatok közül 
szénhidrátban egyik sem javított következetesen. Az előmérésben
a P3\_recipe prompt csökkentette leginkább a zsírbecslés hibáját az OmniFood8k-n, és a döntési mérés 
ezt a javulást a Flash-Lite-nál és a Solnál is szignifikánsan megerősítette, a szénhidrátbecslés 
hibáján viszont egyik adathalmazon sem változtatott.

\paragraph{}A SNAPMe-n a referenciarekordból képzett ételnév-lista a Flash-Lite zsírhibáját 
szignifikánsan csökkentette, a szénhidráthibáját nem változtatta. Mivel a névlista ugyanabból a 
rekordból származik, mint a referenciaértékek, ez felső becslés arra, amit egy felhasználói leírás 
elérhet. A csak szöveges NutriBench-leírásokon a modellek pontossága nem különbözött kimutathatóan, 
a költségük viszont igen.

\paragraph{}A Flash-Lite minden adathalmazon a leggyorsabb és a legolcsóbb modell volt. A döntési 
mérés 5546 hívása között nem volt sikertelen hívás vagy fel nem dolgozható válasz, a JSON-sémának 
megfelelő válaszok között azonban négy nagyságrendi hibát is találtam. Az ugyanarra az ételre adott 
két becslés szénhidrátban az ételenkénti szórás alapján a referenciaátlag 4,6--27,2\%-ával 
ingadozott. A modellek által megnevezett tételek a Gemini 3.8 Flash leképezésével a referencianevek 
22,8--42,1\%-át fedték le, és ez az érték a referencia részletességétől és a leképező modelltől is 
függött.

\paragraph{}Az eredmények alapján egy étkezésnaplózó alkalmazáshoz a Flash-Lite egy képpel és a 
P3\_macros prompttal a vizsgált modellek közül a legkedvezőbb költségű kiindulópont, amely a 
legtöbb vizsgált adathalmazon szénhidrátban a drágább modelleknél is pontosabb volt. A becslés hibája 
azonban ételenként nagy lehet, ezért az eredményt a felhasználónak ellenőrizhető és javítható 
javaslatként érdemes megjeleníteni.

\section{Jövőbeli tervek}
\subsection{Javasolt becslési folyamat}
\label{sec:javasolt_folyamat}
\paragraph{}A javasolt folyamatban a felhasználó lefényképezi az étkezést, az alkalmazás a képet a 
mérésekkel azonos módon előkészíti, vagyis az EXIF-adatok szerint elforgatja, a hosszabbik oldalát 
legfeljebb 1024 képpontra kicsinyíti és JPEG formátumba kódolja, majd egyetlen hívással, JSON-séma 
szerinti válaszkéréssel elküldi a modellnek. A javasolt alapprompt a P3\_macros, mert a cukorbetegség szempontjából 
legfontosabb szénhidrátbecslésen a P3\_recipe nem javított, és a Flash-Lite átlagos válaszidejét növelte. 
Ha a zsírbecslés pontossága is fontos, a P3\_recipe használata indokolt lehet. A választ az alkalmazás a 
hihetőségi szabályokkal ellenőrzi, és a felismert tételeket a becsült tömeggel és tápértékkel együtt 
a felhasználó elé teszi, aki a tételeket és a mennyiségeket a mentés előtt javíthatja. Az alapeset a 
szinkron feldolgozás, amelyben a felhasználó megvárja a választ, ez a Flash-Lite mért 2,1--2,6~s-os 
medián válaszidejével elfogadható várakozást jelent. Lehetőség a háttérben végzett feldolgozás is, 
amelyben a kérés adatbázisba kerül, és az alkalmazás értesíti a felhasználót, ha a becslés elkészült. 
Ez a lassabb modelleknél vagy rossz hálózati kapcsolat esetén lehet előnyös.
Egy kérés feldolgozásának ideje szinte teljes egészében a szolgáltató válaszára való várakozás, 
az alkalmazás szerverén végzett számítás ehhez képest elhanyagolható. Ezért érdemes a kéréseket 
a szerveren nem blokkoló módon, aszinkron hálózati hívásokkal vagy egy kérési sor segítségével 
kezelni, így a szerver egy válaszra várva is fogadhat és továbbíthat újabb kéréseket, és kevés 
erőforrással egyszerre sok felhasználót szolgálhat ki.

\subsection{További vizsgálatok}
\paragraph{}A mérések adathalmazokon készültek, ezért a következő lépés az alkalmazás leendő 
felhasználóinak étkezésein végzett mérés lenne, a felhasználók saját szöveges leírásával, amely a 
referenciarekordból képzett névlistával szemben a valós szöveges bemenet hatását mutatná meg. Érdemes 
lenne megvizsgálni, hogy több hívás becsléseinek átlagolása mennyit csökkent a hibán és az 
ingadozáson, és a konfidenciaintervallumokat a résztvevőnkénti csoportosítás figyelembevételével is 
kiszámítani. A csúcsszintű modellek teljes körű mérése, valamint a képen látható ismert méretű 
tárgyak és a felhasználó által megadott adag felhasználása szintén további vizsgálatot igényel.
